\chapter{Objetivos y requisitos}

\section{Objetivo general}
Desarrollar una aplicación de escritorio que apoye la gestión operativa de una tienda de electrónica mediante módulos diferenciados para administración, punto de venta y almacén, con autenticación por roles, persistencia en archivos CSV e interfaz gráfica JavaFX.

\section{Objetivos específicos}
\begin{enumerate}
\item Validar el acceso de cada empleado contra su registro persistido, con rol explícito.
\item Permitir al administrador abrir y cerrar la caja, condicionando con ello la confirmación de ventas.
\item Administrar empleados, productos y proveedores con validaciones y unicidad de identificadores.
\item Informar ventas del día, entradas acumuladas y productos con existencias críticas, con exportación a archivo de texto.
\item Ofrecer al cajero búsqueda por código, nombre o marca con filtro por categoría, carrito con control de stock, proforma en .txt y venta confirmada con garantía asociada.
\item Registrar entradas manualmente o desde CSV con vista previa antes de persistir.
\item Mantener integridad referencial suficiente entre ventas, clientes y garantías.
\end{enumerate}

\section{Alcance del sistema}
El sistema cubre la autenticación de empleados, la gestión de caja, el catálogo e inventario, las entradas de productos, las ventas y proformas, las garantías, los clientes, los proveedores y los pedidos de proveedor. Incluye exportación de informes de ventas y entradas desde el panel administrativo. La solución no incluye una base de datos relacional, sincronización entre sucursales, un servicio web ni un mecanismo externo de copias de seguridad.

\section{Usuarios y necesidades}
\begin{table}[H]\centering\caption{Usuarios del sistema y necesidades principales}\begin{tabular}{p{.22\textwidth}p{.62\textwidth}}\toprule
Usuario & Necesidades y responsabilidades\\\midrule
Administrador & Controlar la caja, supervisar ventas y entradas, consultar inventario (completo y crítico), administrar empleados, productos y proveedores, exportar informes.\\
Cajero & Buscar productos y filtrar por categoría, preparar carrito con control de stock, emitir proforma, confirmar venta con datos de cliente obligatorios (nombre, CI y teléfono), consultar garantías y anular operaciones.\\
Almacenero & Consultar el catálogo e inventario actual, registrar productos nuevos, incrementar existencias y registrar entradas manuales o importadas desde CSV.\\\bottomrule
\end{tabular}\end{table}

\section{Requisitos funcionales}
\begin{enumerate}
\item \textbf{RF-01 Autenticación:} validación de nombre, contraseña y rol contra \texttt{empleado.csv}; solo el panel del rol autenticado se abre.
\item \textbf{RF-02 Caja:} apertura y cierre por el administrador; sin caja abierta no se confirma ninguna venta.
\item \textbf{RF-03 Inventario:} vista del inventario completo filtrable por categoría y vista de productos con cinco unidades o menos.
\item \textbf{RF-04 Administración:} alta y baja de empleados (con protección del último administrador), productos y proveedores.
\item \textbf{RF-05 Supervisión:} total de ventas del día, historial de entradas con costo acumulado y exportación de informes .txt.
\item \textbf{RF-06 Venta:} búsqueda incremental por código, nombre o marca con filtro de categoría opcional; categorías derivadas del propio \texttt{inventario.csv}.
\item \textbf{RF-07 Confirmación:} al confirmar, descuenta existencias, registra o actualiza el cliente (identificado por CI), genera garantías de 12 meses y deja constancia en \texttt{ventas.csv} y \texttt{garantias.csv}.
\item \textbf{RF-08 Post-venta:} anulación de una venta o proforma (restituye stock si estaba confirmada), consulta de garantías por código o identificador y exportación del comprobante.
\item \textbf{RF-09 Almacén:} registro de productos nuevos y aumento de existencias, con constancia en \texttt{entradas.csv}.
\item \textbf{RF-10 Importación:} vista previa de entradas u órdenes CSV antes de confirmar la importación.
\end{enumerate}

\section{Requisitos no funcionales}
\begin{itemize}
\item \textbf{RNF-01 Plataforma:} Java 17, JavaFX 17.0.2 y Maven; multiplataforma (Linux y Windows).
\item \textbf{RNF-02 Usabilidad:} ventana única con cambio de sección por pestañas; avisos y formularios en cuadros internos; tema claro/oscuro persistente en la sesión.
\item \textbf{RNF-03 Persistencia:} CSV en UTF-8 con encabezado y entrecomillado cuando el campo lo requiere.
\item \textbf{RNF-04 Integridad:} validación de obligatoriedad, formato de CI y teléfono, stock no excedido, unicidad de CI de empleados, CI de clientes y existencia de identificadores antes de persistir.
\item \textbf{RNF-05 Mantenibilidad:} separación de vistas FXML, controladores, dominio (\texttt{user}), modelos y servicios; nombres de archivos centralizados en \texttt{config.ArchivosCSV}.
\item \textbf{RNF-06 Verificabilidad:} suite JUnit 5 (516 pruebas) ejecutable con \texttt{mvn test}, con casos pendientes documentados en \texttt{src/test/java/TESTS\_PENDIENTES.md}.
\end{itemize}

\section{Tecnologías y organización}
\begin{table}[H]\centering\caption{Tecnologías verificadas}\begin{tabular}{ll}\toprule Tecnología & Evidencia\\\midrule
Java 17 & \texttt{maven.compiler.source/target}\\
JavaFX 17.0.2 & \texttt{javafx-controls}, \texttt{javafx-fxml}\\
Maven & \texttt{javafx-maven-plugin} (\texttt{com.overbyte.Main})\\
JUnit 5.10.0 & \texttt{junit-jupiter} + \texttt{maven-surefire-plugin}\\
CSV & \texttt{MotorCSV} con parseo propio y \texttt{data/} como almacén\\\bottomrule\end{tabular}\end{table}
